\chapter{Computational Environment}
\label{app:computational_environment}

\section{Containerized Execution}
\label{sec:appendix_containerized_execution}

The two tokenizer implementations have incompatible software requirements and
were therefore executed in separate Docker images on a CUDA-enabled server.
The ImageNet datasets and pretrained checkpoints were mounted into the
containers, while experiment directories were mounted into persistent storage.
Long-running full-dataset exports were launched in background terminal
sessions. Random seeds, perturbation modes, patch sizes, and output locations
were fixed by script arguments or environment variables and recorded in run
directory names and metadata.

\section{Tokenizer Environments}
\label{sec:appendix_tokenizer_environments}

The VQGAN environment is based on Ubuntu 20.04 with CUDA 11.7, Python 3.8,
PyTorch 1.13.1, and torchvision 0.14.1. It retains the older dependencies
required by Taming Transformers, including PyTorch Lightning 1.0.8,
OmegaConf 2.0.0, and scikit-image 0.18.3.

The LlamaGen environment uses Ubuntu 20.04 with CUDA 12.2, PyTorch 2.1.2, and
torchvision 0.16.2. Its analysis dependencies include NumPy, Pillow,
scikit-image, pandas, Matplotlib, and LPIPS. Both environments expose the same
dataset and checkpoint mounts, allowing the model-specific exporters to
produce compatible image and metadata layouts.

The model exporters do not use byte-identical resizing operators. LlamaGen
repeatedly downsamples large images with Pillow BOX filtering, resizes the
shortest side with bicubic interpolation, and then applies an integer center
crop. The VQGAN exporter resizes the shortest side with Lanczos interpolation
before center cropping. Both produce RGB $256\times256$ tensors, but this
filtering difference is a preprocessing confound in fine-grained cross-model
comparisons.

\section{Evaluation and Analysis Environment}
\label{sec:appendix_evaluator_environment}

Distribution-level reconstruction metrics are computed in a separate
evaluator image containing TensorFlow CPU 2.9, PyTorch, torchvision, and
\texttt{torch-fidelity}. Reconstructions are exported as uint8 NPZ arrays and
compared with reference statistics generated from the same ImageNet validation
preprocessing. Using a separate evaluator avoids changing either tokenizer
environment solely to satisfy metric dependencies.

The chapter notebooks run in the Jupyter environment associated with
\texttt{vq-explore}. They use pandas and NumPy for aggregation, Matplotlib for
figures, scikit-image for pixel metrics, and batched LPIPS evaluation where
required. Full ImageNet experiments use 50{,}000 validation images. DataLoader
workers, asynchronous image writers, process pools, and GPU metric batches are
used to keep model inference and analysis practical without changing the
experimental definitions.

The exact GPU model and wall-clock duration of every historical run were not
archived consistently and are therefore not reconstructed here. Container
versions, checkpoints, seeds, scripts, and compact outputs are retained; the
omitted hardware timing details do not affect the numerical definitions of the
reported metrics.
